\section{Introduction}\label{sec:intro}

Let $\mathcal T_n$ be the set of domino tilings of a $3\times 2n$
rectangle, where $n\ge1$. A domino covers two unit cells sharing an
edge. Both orientations are allowed, the dominoes are unlabelled,
and no tatami restriction is imposed. Thus a tiling is a set of
unordered pairs of adjacent cells.
Number rows from top to bottom and columns from left to right.
The maps
\[
 h(r,c)=(4-r,c),\qquad v(r,c)=(r,2n+1-c),\qquad \rho=hv
\]
are reflection in the horizontal axis, reflection in the vertical
axis, and rotation through $180^\circ$, respectively.
Since $3\ne2n$, the full symmetry group of the rectangle is
$G=\{e,h,v,\rho\}\cong C_2\times C_2$.
Write $a(n)=|\mathcal T_n/G|$. We set $\mathcal T_0$ to consist of a
single empty tiling, fixed by all four abstract group elements, so
that $a(0)=1$.

Mathar \cite[Table 10, Conjecture 6]{mathar2013} proposed a generating
function for these symmetry classes. The equivalent formula in
OEIS A232622 \cite{oeis} is accompanied by a conjectured recurrence
of order seven. We present two candidate derivations for further verification.
The first uses a complete classification of vertically
indecomposable blocks and the action of $G$ on words in these blocks.
The second counts symmetric tilings directly from partial tilings
of a half-rectangle.
The latter construction also yields a uniform formula for any fixed
odd height.

Transfer matrices and indecomposable tilings are established tools
in fixed-width enumeration. In particular, Mathar
\cite[Sections 2.1--2.3]{mathar2014} gives the ordinary three-row
generating function and the corresponding indecomposable-block
generating function. We rederive the required block classification
and determine the symmetry action explicitly.
DeFord \cite[Theorem 6.1]{deford} establishes a general recurrence
framework for rectangular tilings counted up to symmetry and
illustrates it for two-row domino tilings in Section 6.2.1.
Thus rationality at fixed height is established background;
our odd-height formula specifies the boundary conditions and
gives an explicit five-row evaluation.
None of our
combinatorial arguments assumes the target recurrence or a finite
list of coefficients.
The source entry was checked on October 25, 2026; its conjecture
label does not establish that no equivalent proof appears elsewhere.
No claim of bibliographic priority is made here.

\subsection{Main result}
Let
\[
 d_n=|\mathcal T_n|,\qquad h_n=|\Fix(h)|,\qquad
 v_n=|\Fix(v)|,\qquad r_n=|\Fix(\rho)|.
\]
Their ordinary generating functions, with $x$ marking the
half-width $n$, are denoted by $D,H,V,R$.
All generating-function identities are interpreted in
$\mathbb Q[[x]]$.

\begin{theorem}\label{thm:main}
The number of domino tilings of a $3\times2n$ rectangle up to
rectangle symmetries has generating function
\begin{equation}\label{eq:main}
 A(x):=\sum_{n\ge0}a(n)x^n
 =\frac{1-3x-4x^2+10x^3-4x^5+x^6}
 {(1-x)(1-4x+x^2)(1-4x^2+x^4)}.
\end{equation}
Consequently,
\begin{equation}\label{eq:recurrence}
\begin{split}
a(n)={}&5a(n-1)-a(n-2)-19a(n-3)+19a(n-4)\\
       &+a(n-5)-5a(n-6)+a(n-7),\qquad n\ge7,
\end{split}
\end{equation}
with initial values
\[
 (a(0),\ldots,a(6))=(1,2,5,14,46,156,561).
\]
\end{theorem}

Multiplying the numerator and denominator in \eqref{eq:main} by
$-1$ gives the form displayed in A232622.
Sections \ref{sec:blocks} and \ref{sec:half} give the two proofs.
Section \ref{sec:types} determines the exact stabilizers and the
asymptotic proportion of symmetric tilings.
Sections \ref{sec:odd} and \ref{sec:five} treat arbitrary odd height
and the explicit five-row case.

\begin{lemma}[Burnside's formula]\label{lem:burnside}
For every $n\ge0$,
\begin{equation}\label{eq:burnside}
 a(n)=\frac{d_n+h_n+v_n+r_n}{4}.
\end{equation}
\end{lemma}
\begin{proof}
Count the pairs $(g,T)\in G\times\mathcal T_n$ with $gT=T$:
\[
 \sum_{g\in G}|\Fix(g)|
 =\sum_{T\in\mathcal T_n}|\Stab(T)|.
\]
By the orbit--stabilizer formula, each orbit $O$ contributes
$|O|\cdot4/|O|=4$ to the right-hand side.
Division by four proves the assertion.
This is the usual orbit-counting lemma; see
\cite[Chapter 7, Lemma 7.2]{stanley}.
The argument does not require a faithful action and includes $n=0$.
\end{proof}
